\section*{Recap of Lecture 4}
\phantomsection
\makeatletter
\def\@currentlabel{Recap of Lecture~4}\def\@currentlabelname{Recap of Lecture~4}%
\def\cref@currentlabel{[recap][1][]Recap of Lecture~4}%
\makeatother
\crefformat{recap}{#2#1#3}\Crefformat{recap}{#2#1#3}
\label{sec:recap}
\addcontentsline{toc}{section}{Recap of Lecture 4}

Lecture~4 opened the tools node of \msa{}, the node of the course's running agent that runs the
tools the model requests. It studied the harness, the program that runs the model loop and executes
tool calls for the model. This recap restates the parts of that lecture on which this lecture builds: the
tool contract, the dispatcher, the guarantees of the harness, and the integration problem that the
Model Context Protocol (MCP) addresses. The Lecture~4 notes, \emph{Tool calls, contracts and
execution boundaries}, give the full treatment.

A tool contract is the written specification of a tool, and Lecture~4 divided it into five parts:
\bi
\item \textbf{Purpose.} The contract states what the tool does.
\item \textbf{Arguments.} The contract lists the arguments that the tool accepts, with the
constraints on their values.
\item \textbf{Results and errors.} The contract states what a successful call returns and how a
failed call is reported.
\item \textbf{Side effects.} The contract states the changes a call makes besides returning its
result, such as changes to files.
\item \textbf{Execution constraints.} The contract limits where and how the tool may operate, for
example by requiring that every file the tool writes lies inside the run's own copy of the code
base.
\ei
\msa{} has three tools. \demoref{read\_file} and \demoref{search} only read, and \demoref{edit} is
the only tool that writes files.

The model receives only part of the contract, the tool definition. It consists of the tool's name,
a description of its purpose, and the argument schema, which gives the names and types of the
arguments and the constraints on their values. The tool's implementation and the dispatcher, the
component of the harness that executes each tool call, define the rest of the contract: the
results, the errors, the side effects and the execution constraints.

A description is text that the model reads. It can state a rule, but nothing in it stops a call
that breaks the rule, so a description alone enforces nothing. A constraint holds only where code in
the implementation or in the dispatcher checks it. \Cref{sec:mcpcontract} returns to the five parts
of the contract and asks which of them MCP enforces, and \cref{sec:attacks} uses the same parts to
describe what an attacker alters in a tool.

The workspace is the disposable copy of the code base that one run may read and edit. The
dispatcher receives each tool request, executes it against the workspace, and returns an
observation to the model. An observation is the tool's result as the harness records it.

The dispatcher holds the registry, the workspace root and the trace, and the model never touches any
of them. The registry maps each tool name to the function that implements the tool. The workspace
root is the path of the workspace, and the trace is the log of events that the harness writes
during a run.

The tools node passes each tool call in the model's reply to the dispatcher. The dispatcher finds
the tool in the registry, binds the model's arguments to the parameters of the tool function, and
calls the function with the workspace root, which it supplies itself. It records the request and its
result in the trace and returns the observation, which the tools node turns into a tool message, the
message that carries one tool result back to the model. The model never calls a tool
function directly. It can only ask for a call, and the dispatcher decides whether and how the call
runs.

In \msa{}, the dispatcher, the tool functions and their checks all run in the agent's own process.
\Cref{sec:mcp} begins by explaining why a check that runs inside that process can be switched off or
avoided.

A guarantee is a property that the harness promises for every run, such as that only registered
tools run. Lecture~4 asked which component enforces each guarantee, and
\cref{tab:recap-guarantees} collects the answers. The row to look at is the one about reads, the
only guarantee that no component enforces.

\begin{table}[htbp]
  \caption{The guarantees of the harness of \msa{} and where each is enforced. Reads are not
  confined to the workspace.}
  \label{tab:recap-guarantees}
  \begin{adjustbox}{width=\linewidth}\begin{tabularx}{\linewidth}{@{}L{0.36\linewidth}X@{}}
    \toprule
    \textbf{Guarantee} & \textbf{Where it is enforced} \\
    \midrule
    \textbf{Only registered tools run} & The dispatcher checks the registry before calling a tool. \\
    \textbf{Arguments match the tool's parameters} & The dispatcher binds the arguments and converts
      unambiguous type mismatches. \\
    \textbf{Writes stay inside the workspace} & The \demoref{edit} tool, the only writer, checks
      paths with the path check. \\
    \textbf{Reads stay inside the workspace} & Not enforced. The glob of \demoref{search} bypasses
      the path check. \\
    \textbf{No model call after the budget is exhausted} & Each model node checks the budget before
      calling the model. \\
    \bottomrule
  \end{tabularx}\end{adjustbox}
\end{table}

The path check, through which \demoref{edit} enforces the guarantee about writes, accepts a path
only if it resolves to a location inside the workspace root. The model nodes are the nodes that
call the model. Each of them checks the budget, the limit on the number of model calls in a run,
before it calls the model. Every enforcer in the table is code that runs before the effect it protects.

Reads are not confined, because \demoref{search} skips the path check. The model may supply
\demoref{search} with a glob, a file pattern with wildcards that selects the files to search.
\demoref{search} applies the glob to the workspace root and opens every file that the glob matches,
so a glob that climbs out of the root lets the tool read files outside the workspace. The path
check itself is correct. \demoref{search} is a route to the files that does not pass through it.

Native tool calling, studied in Lecture~4, makes the tool call part of the model API, with one
format per vendor. It standardizes the exchange between a model and an application, and it leaves
the integration between an application and a tool to each application. An integration is the code
that lets one application use one tool: the tool's definition in the application's format, and the
code that carries out the five operations of a tool call. That code finds the tool, binds the
arguments, checks policy (the rules that decide whether a requested call is allowed), runs the tool
and observes the result. The work multiplies for the following reasons:
\bi
\item \textbf{Vendor formats.} Each vendor defines its own format, so every application converts
between formats. Each vendor's strict mode, which makes the model's arguments match the argument
schema, accepts a different subset of JSON Schema, so even the argument schema must be adapted for
each vendor.
\item \textbf{Repeated operations.} Each application re-implements the five operations for every
tool it uses.
\item \textbf{Place and language.} Tools live in the process and language of the application that
uses them, so a tool written as a Python function cannot serve an application written in
TypeScript.
\ei

As a result, $N$ applications and $M$ tools need $N \times M$ integrations, one for each pair, which
is the $N \times M$ integration problem. A shared protocol reduces the count to $N + M$, because
each application and each tool implements the protocol once. For three applications and three
tools, the count falls from nine to six. MCP, which Lecture~4 introduced
briefly, is such a protocol for the tool side, and \cref{sec:mcp} studies it in full.

This lecture reads \cref{tab:recap-guarantees} from the attacker's side, where each unenforced row
marks a fault that an adversary can use (\cref{sec:threats}), and it moves the tools out of the
harness.
